\section{Problemstellung und Motivation}
\label{sec:problemstellung}

Die globale Versicherungsbranche befindet sich in einer Phase tiefgreifender
digitaler Transformation. Steigender Wettbewerbsdruck, verändertes
Kundenverhalten und eine rasant zunehmende Menge strukturierter und
unstrukturierter Daten zwingen Versicherungsunternehmen dazu, ihre
analytischen Fähigkeiten grundlegend zu modernisieren
\citep{eiopa2023digitalisation}. Neben klassischen, auf historisierten Daten
beruhenden Analysen gewinnt daher der Einsatz moderner
Data-Science-Methoden an Bedeutung. Das Anwendungsspektrum reicht von der
KI-gestützten Betrugserkennung im Schadenfall bis zu prädiktiven
Retention-Modellen im Kundenbeziehungsmanagement
\citep{debener2023,spiteri2018}.

Als weltweit agierender Konzern steht die Allianz SE vor der Herausforderung,
große Datenmengen aus heterogenen Quellen effizient zu konsolidieren und
analytisch nutzbar zu machen. Zu diesen Quellen zählen unter anderem
Bestands-, Finanz- und Kundeninteraktionsdaten sowie Schadensfotos und
Telemetriedaten. Die zentrale Problemstellung besteht darin, eine flexible und
skalierbare Datenarchitektur zu konzipieren, die sowohl performante operative
Abfragen mit geringen Latenzen als auch stichtagsbezogene
Business-Intelligence-Auswertungen und das Training von
Machine-Learning-Modellen ermöglicht \citep{schulz2024}.
